\section{Konzeption der Migrations-Pipeline} \label{konzeption}

Dieses Kapitel beschreibt das konzeptionelle Design der orchestrierten Migrations-Pipeline.

\subsection{Rollen im orchestrierten Workflow}

Der Workflow wird in spezialisierte Rollen zerlegt, um nicht-deterministische Modellantworten durch klare Zuständigkeiten und Zwischenschritte zu kontrollieren. Dadurch entsteht eine Pipeline, in der Ergebnisse schrittweise weiterverarbeitet und überprüft werden. Vergleichbare agentische Rollenkonzepte wurden etwa in ChatDev für Softwareentwicklungsprozesse untersucht \parencite{qian2024chatdev}.

\subsubsection{Analyzer}

Der Analyzer erschließt die bestehende Codebasis mithilfe statischer Analyse und strukturierter Code-Intelligence. Ziel ist die Gewinnung fachlicher und technischer Informationen auf Basis von Symbol- und Abhängigkeitsbeziehungen.

\subsubsection{Spec-Writer}

Der Spec-Writer überführt die Analyseergebnisse in textuelle Spezifikationen. Diese Markdown-Artefakte sollen sowohl für menschliche Prüfung als auch für nachgelagerte Generierungsschritte nutzbar sein.

\subsubsection{Verifier}

Der Verifier prüft, ob die Spezifikation mit dem rekonstruierbaren Quellkontext übereinstimmt. Dadurch wird verhindert, dass implizite Annahmen oder halluzinierte Details unbemerkt in die Implementierung einfließen.

\subsubsection{Builder}

Der Builder erzeugt aus der validierten Spezifikation den Zielcode. Die Erzeugung basiert damit nicht direkt auf dem Ausgangscode, sondern auf einem geprüften Zwischenartefakt.

\subsection{Feedback und iterative Validierung}

Die Pipeline ist als iterativer Prozess mit Rückkopplungsschleifen zu verstehen. Analyseergebnisse, Spezifikationen und Verifikationsergebnisse beeinflussen sich wechselseitig, sodass unklare oder fehlerhafte Ableitungen frühzeitig korrigiert werden können.

\subsubsection{Feedback-Schleifen zwischen Spezifikation und Quell-Kontext}

Feedback-Schleifen zwischen Spezifikation und Quell-Kontext ermöglichen eine schrittweise Schärfung der Anforderungen. Erkennt der Verifier Unstimmigkeiten, müssen Analyzer und Spec-Writer ihre Artefakte gezielt nacharbeiten, bevor eine Implementierung zulässig ist.

\subsubsection{Abgleich mit dem Code Knowledge Graph}

Der automatische Abgleich gegen den Code Knowledge Graph liefert dafür eine deterministische Referenz. Er ermöglicht den Vergleich modellgenerierter Aussagen mit Symbolen, Abhängigkeiten und Strukturen der Ausgangsbasis.

\subsection{Phasenmodell der Pipeline}

Das Phasenmodell der Pipeline gliedert sich in Ingestion, Spec-Generation und Code-Synthese. In der Ingestion wird die Codebasis strukturiert erschlossen, in der Spec-Generation in ein überprüfbares Zwischenartefakt überführt und in der Code-Synthese in eine Zielimplementierung transformiert.
